\scanpage{19}
\mspage{9}
On sait \inserted{alors} que l'on a la notion d'application $X\longrightarrow Y$ entre espaces modérés, qui est \uline{modérée au voisinage d'un pt}, et on trouve

\uline{Prop} Pour qu'une application d'un espace modéré $X$ dans un espace ind-modéré soit modérée, il faut s. qu'elle le soit au voisinage de chaque pt.

Donc les applications modérées de $X$ dans $Y$ forment les sections d'un faisceau --- le faisceau des applications modérées de $X$ dans $Y$. La structure \inserted{modérée} de $X$ est connue quand on connaît l'ens des applications modérées de $X$ dans $R$ (cela \unsure{vaut|faut} dans le contexte purement ensembliste, avec $(M1)_0$\unsure{;|?} $(M2)$ $(M3)$ $(M4)$ seulement), donc aussi quand on connaît le faisceau des germes de fonctions modérées.
\mspage{18}
